\section{Study Design}

\subsection{Overall Design}
\draftinstr{Follow the Phase 2 publication section
\texttt{final-protocol/publication/sections/sec-04-design.tex} and the Phase 1
model file \texttt{trial-protocol/draft-protocol/sections/sec-04-design.tex}.
Describe the prospective, multicenter, randomized, parallel-group, controlled,
open-label design with BICR across eight academic HPB centers; $n=220$ randomized
1:1 (110 per arm); Arm A daraxonrasib RP2D 300 mg PO daily plus PancreSpeed II,
Arm B modified FOLFIRINOX \cite{Conroy2018} plus standard Whipple; central
permuted-block randomization stratified by resectability, KRAS allele, neoadjuvant
therapy, and site. Source the design and population order from
\texttt{trial-protocol/nih-protocol/03\_study\_design\_and\_study\_population.md}.}

\draftinstr{Render \texttt{trial-phase-2/mermaid/fig-08-randomization-design.md} as
a TikZ \texttt{mermaidfig} (this is \textbf{Figure 8}): eligible participant ->
central stratified randomization 1:1 -> Arm A / Arm B -> eight harmonized HPB sites
-> BICR plus central masked pathology -> endpoint adjudication.}

\subsection{Scientific Rationale for Study Design}
\draftinstr{Justify the randomized controlled design as the ethical next step that
inverts the single-arm predicate \cite{kawchak2026phase1}; open-label conduct is
unavoidable for a surgical-system intervention, so bias control is concentrated
downstream (BICR imaging, central masked R0/MPR pathology, blinded adjudication).
The internal control Arm B replaces the external Dutch benchmark
\cite{DutchCohort2025,Siegel2025}. Follow the CONSORT framework
\cite{Schulz2010consort}.}

\subsection{Justification for Dose and Device Readiness}
\draftinstr{Daraxonrasib at the established RP2D of 300 mg once daily (no 3+3
escalation), anchored to RASolute 302 \cite{rasolute302,kawchak2026phase1}. The
device ``dose'' analogue is the \emph{upgraded} Phase 0 gate: USL $\geq 8.0$, at
least 5000 sims across at least 3 frameworks, sim-to-real $< 1.5$ mm and $< 0.4$ N,
per-site qualification plus fleet harmonization, operating with the intra-op hard
caps ($\leq 3$ N per-arm, $\leq 18$ N cumulative, $\leq 3$ ms e-stop, five-vessel
no-fly gate) \cite{ecfr2026part812,kawchak2026adapt312}. Source from
\texttt{trial-protocol/inputs/21cfr312\_adapt/02\_ind\_content\_phases.tex}.}

\draftinstr{Render \texttt{trial-phase-2/mermaid/fig-09-staged-autonomy.md} as a
TikZ \texttt{mermaidfig} (this is \textbf{Figure 9}): the phase-graduated
staged-autonomy model running at Stage 2 supervised semiautonomy (Class II,
continuous oversight, immediate override); any expanded subtask requires
independent safety review and USL $\geq 8.0$; Stage 3 full autonomy reserved and
prohibited under 21 CFR \S312.21(e).}

\subsection{End of Study Definition}
\draftinstr{Define individual completion at the 24-month OS assessment (or death,
progression precluding follow-up, or withdrawal), and overall end of study as the
last scheduled PFS/OS assessment plus primary-analysis database lock after $\sim
140$ PFS events. Discontinued participants (\S7) remain in long-term survival
follow-up.}
